\chapter{Filesystem Layers}
\label{ch:fslayers}

Veilbox implements a multi-layered filesystem strategy that separates
the immutable root filesystem from persistent state storage. This chapter
documents each filesystem layer and its role.

\section{Storage Architecture}

Veilbox uses three distinct storage tiers:

\begin{figure}[H]
\centering
\begin{tikzpicture}[node distance=0.8cm]
\node[highlight] (initramfs) {Initramfs (tmpfs)};
\node[box, below of=initramfs] (overlay) {OverlayFS (container layers)};
\node[box, below of=overlay] (state) {State Partition (ext4)};
\node[box, below of=state] (disk) {Disk Image (raw file)};
\draw[arrow] (initramfs) -- node[right] {read-only /} (overlay);
\draw[arrow] (overlay) -- node[right] {/mnt/state} (state);
\draw[arrow] (state) -- node[right] {/dev/vda} (disk);
\end{tikzpicture}
\caption{Filesystem layering. The initramfs root is read-only at runtime;
containers use OverlayFS on top; persistent data lives on the ext4
state partition.}
\label{fig:fslayers}
\end{figure}

\section{Layer 1: Initramfs (tmpfs)}

The root filesystem is a tmpfs created by the kernel from the embedded
cpio archive:

\begin{itemize}[nosep]
  \item \textbf{Type}: tmpfs (memory-backed, no block device).
  \item \textbf{Contents}: BusyBox, containerd, runc, nerdctl, Dropbear,
        configuration files.
  \item \textbf{Size}: Approximately 90 MB uncompressed.
  - \textbf{Persistence}: None — every reboot gets a fresh copy from the
        kernel binary.
  - \textbf{Mount}: \texttt{mount -t tmpfs none /}
\end{itemize}

\section{Layer 2: State Partition (ext4)}

The state partition provides persistent storage across reboots:

\begin{itemize}[nosep]
  \item \textbf{Location}: \texttt{/mnt/state}
  - \textbf{Device}: \texttt{/dev/vdb} (external VirtIO disk) or
        \texttt{/dev/sdb} (external SATA disk) or \texttt{/state/} on
        the VEILBOX partition (fallback).
  - \textbf{Filesystem}: ext4.
  - \textbf{Size}: 128 MB (default).
  - \textbf{Contents}: containerd data (images, snapshots, metadata),
        container logs, SSH host keys, /var/log/messages.
\end{itemize}

\subsection{Mount Strategy}

The rcS script uses a fallback strategy:

\begin{lstlisting}[language=bash]
if [ -b /dev/vdb ]; then
    mount /dev/vdb /mnt/state
elif [ -b /dev/sdb ]; then
    mount /dev/sdb /mnt/state
elif [ -d /state ]; then
    mount --bind /state /mnt/state
else
    mount -t tmpfs none /mnt/state
fi
\end{lstlisting}

\section{Layer 3: OverlayFS (Container Layers)}

Containers use OverlayFS for copy-on-write filesystem layering:

\begin{itemize}[nosep]
  \item \textbf{Lower dir}: Image layers (read-only snapshots).
  \item \textbf{Upper dir}: Container writable layer (changes since
        container creation).
  \item \textbf{Merge dir}: The unified view presented to the container
        process.
\end{itemize}

\section{Disk Image Structure}

The 8 GB raw disk image (\texttt{veilbox.raw}) uses MBR partitioning:

\begin{lstlisting}
Disk veilbox.raw: 8 GiB, 8589934592 bytes, 16777216 sectors
Device        Boot Start     End Sectors Size Id Type
veilbox.raw1  *     2048 4196351 4194304   2G 83 Linux
\end{lstlisting}

The partition is formatted as ext4 with the label \texttt{VEILBOX}.


\section{Layer 4: Container Image Layers}

Container images are composed of OCI-compliant layers. Each layer is a
tar archive of filesystem changes:

\begin{lstlisting}
# Example: nginx image layers
Layer 0 (base): alpine:latest  ->  /bin/sh, /lib, /etc
Layer 1:        nginx install  ->  /usr/sbin/nginx, /etc/nginx
Layer 2:        config change  ->  /etc/nginx/nginx.conf
Layer 3 (writable): container  ->  /var/log/nginx/access.log
\end{lstlisting}

Each layer is stored as a gzip-compressed tar archive in the containerd
content store. The digest (SHA256) uniquely identifies each layer.

\section{Filesystem Hierarchy Standard (FHS)}

Veilbox implements a subset of the Filesystem Hierarchy Standard:

\begin{longtable}{|p{3cm}|p{4cm}|p{5cm}|}
\hline
\textbf{Directory} & \textbf{Purpose} & \textbf{Type} \\
\hline
/bin & Essential user binaries & initramfs \\
/sbin & System binaries & initramfs \\
/etc & Configuration files & initramfs \\
/usr/bin & Non-essential binaries & initramfs \\
/usr/sbin & Non-essential system bins & initramfs \\
/lib64 & Shared libraries & initramfs \\
/proc & Process information & procfs \\
/sys & Kernel/device information & sysfs \\
/dev & Device nodes & devtmpfs \\
/tmp & Temporary files & tmpfs \\
/var/log & Log files & state (bind mount) \\
/run & Runtime data & tmpfs \\
/mnt/state & Persistent storage & ext4 (external) \\
/root & Root user home & initramfs \\
\hline
\end{longtable}

\section{Debugfs Rootless File Operations}

The build script uses \texttt{debugfs} to write files to the ext4
filesystem without mounting it. This avoids the need for loop device
access:

\begin{lstlisting}[language=bash]
# Write a file to the ext4 image
debugfs -w output/veilbox.raw <<EOF
mkdir /boot
write /path/to/vmlinuz /boot/vmlinuz
write /path/to/grub.cfg /boot/grub/grub.cfg
EOF
\end{lstlisting}

Key debugfs commands:
\begin{itemize}[nosep]
  \item \texttt{write <host> <target>}: Copy host file to image.
  \item \texttt{mkdir <path>}: Create directory.
  \item \texttt{rm <path>}: Remove file.
  \item \texttt{stat <path>}: Show file metadata.
  \item \texttt{ls -l <path>}: List directory.
  \item \texttt{cd <path>}: Change directory.
\end{itemize}
